\documentclass[a4paper,12pt]{article}

% --- Packages ---
\usepackage[utf8]{inputenc}
\usepackage[T1]{fontenc}
\usepackage[polish]{babel}
\usepackage{geometry}
\geometry{margin=2.5cm}
\usepackage{xcolor}
\usepackage{titlesec}
\usepackage{hyperref}
\usepackage{enumitem}
\usepackage{tcolorbox}
\usepackage{graphicx}
\usepackage{lastpage}
\usepackage{fancyhdr}
\usepackage{seqsplit}
\usepackage{longtable}

% --- Colors ---
\definecolor{primary}{HTML}{2563EB}
\definecolor{secondary}{HTML}{64748B}
\definecolor{success}{HTML}{22C55E}
\definecolor{warning}{HTML}{F59E0B}
\definecolor{danger}{HTML}{EF4444}
\definecolor{bg}{HTML}{F8FAFC}

% --- Header & Footer ---
\pagestyle{fancy}
\fancyhf{}
\setlength{\headheight}{15pt}
\rhead{\textcolor{secondary}{Projekt Terminal - Raport z Audytu Technicznego}}
\lhead{\textcolor{secondary}{Czerwiec 2026}}
\rfoot{Strona \thepage\ z \pageref{LastPage}}

% --- Custom Commands ---
\newcommand{\finding}[4]{
    \begin{tcolorbox}[colframe=#1, colback=white, title={#2}, fonttitle=\bfseries]
        \textbf{Lokalizacja:} \texttt{\seqsplit{#3}} \\
        \textbf{Opis:} #4
    \end{tcolorbox}
}

\newcommand{\successbox}[2]{
    \begin{tcolorbox}[colframe=success, colback=white, title={#1}, fonttitle=\bfseries]
        \textbf{Status:} #2
    \end{tcolorbox}
}

% --- Document Info ---
\title{
    \vspace{2cm}
    \textcolor{primary}{\Huge\textbf{Kompleksowy Raport z Audytu Systemu TERMINAL}}\\
    \vspace{1cm}
    \large Wersja 2.0 - Rozszerzona
}
\author{Zespół Audytu Gemini CLI}
\date{13 Czerwca 2026}

\begin{document}

\maketitle
\thispagestyle{empty}
\newpage

\tableofcontents
\newpage

\section{Podsumowanie Wykonawcze}
Niniejszy dokument stanowi szczegółowy raport z audytu technicznego systemu „Terminal”, przeprowadzonego w czerwcu 2026 roku. Audyt objął analizę kodu źródłowego (frontend i backend), konfigurację infrastruktury, bezpieczeństwo danych oraz procesy CI/CD.

\subsection{Kluczowe Wnioski}
Mimo solidnych fundamentów architektonicznych opartych na Czystej Architekturze (Clean Architecture) i nowoczesnym stosie technologicznym (\textbf{.NET 9, React 18, PostgreSQL}), system wykazuje krytyczne luki w obszarze bezpieczeństwa operacyjnego oraz drastyczne braki w pokryciu testami automatycznymi.

\begin{itemize}
    \item \textbf{Ogólna Ocena:} \textcolor{danger}{\textbf{C+}}
    \item \textbf{Status Bezpieczeństwa:} Krytyczny (wymaga natychmiastowej interwencji).
    \item \textbf{Dług Techniczny:} Średni, skoncentrowany w warstwie infrastruktury i brakujących optymalizacji frontendu.
\end{itemize}

\section{Metodologia Audytu}
Audyt został przeprowadzony z wykorzystaniem następujących technik:
\begin{enumerate}
    \item \textbf{Statyczna Analiza Kodu (SAST):} Manualny i automatyczny przegląd kodu pod kątem wzorców podatności.
    \item \textbf{Analiza Konfiguracji:} Przegląd plików Docker, Ansible oraz konfiguracji CI/CD w GitHub Actions.
    \item \textbf{Przegląd Zależności:} Weryfikacja wersji bibliotek pod kątem znanych podatności (CVE) oraz cyklu życia oprogramowania.
    \item \textbf{Analiza Architektoniczna:} Ocena separacji warstw i przestrzegania zasad SOLID.
\end{enumerate}

\newpage
\section{Szczegółowa Analiza Bezpieczeństwa}

\subsection{Uwierzytelnianie i Autoryzacja}
System wykorzystuje standard JWT z algorytmem HS256. Choć sama implementacja generowania tokenów jest poprawna, zidentyfikowano braki w mechanizmach obronnych API.

\finding{danger}{Brak Rate Limitingu}{Server/src/Terminal.Backend.Api/Program.cs}{
    API nie ogranicza liczby zapytań na sekundę. W dobie zautomatyzowanych ataków, brak limitów (Rate Limiting) naraża infrastrukturę na paraliż (DoS) oraz ułatwia ataki słownikowe na endpoint \texttt{/api/users/login}.
}

\finding{warning}{Podatność IDOR (Insecure Direct Object Reference)}{Server/src/Terminal.Backend.Api/Modules/}{
    Mechanizmy autoryzacji sprawdzają jedynie „czy użytkownik ma uprawnienie X”, zamiast „czy użytkownik ma uprawnienie X do zasobu Y”. Napastnik z uprawnieniami \texttt{ProjectDelete} może teoretycznie usunąć dowolny projekt, znając jedynie jego ID (GUID), nawet jeśli nie jest jego właścicielem.
}

\subsection{Bezpieczeństwo Sieciowe i Nagłówki}
Konfiguracja serwera HTTP w .NET nie wykorzystuje nowoczesnych mechanizmów zabezpieczających przeglądarkę.

\finding{danger}{Brak Nagłówków Security}{Server/src/Terminal.Backend.Infrastructure/Extensions.cs}{
    Brak konfiguracji \texttt{HSTS} (wymuszanie HTTPS), \texttt{Content-Security-Policy} (ochrona przed XSS) oraz \texttt{X-Content-Type-Options}. To sprawia, że użytkownicy końcowi są podatni na ataki typu Man-in-the-Middle oraz wstrzykiwanie skryptów.
}

\finding{danger}{Zbyt Liberalny CORS}{Server/src/Terminal.Backend.Infrastructure/Extensions.cs}{
    Polityka \texttt{AllowAnyOrigin()} na produkcji to „zaproszenie” dla ataków CSRF. API powinno akceptować żądania jedynie z zaufanej domeny frontendu.
}

\newpage
\section{Infrastruktura i Łańcuch Dostaw}

\subsection{Konteneryzacja (Docker)}
\finding{warning}{Praca na Uprawnieniach Root}{Client/Dockerfile, Server/src/Dockerfile}{
    Oba pliki Dockerfile nie definiują użytkownika \texttt{USER}. Procesy wewnątrz kontenera działają z pełnymi uprawnieniami root, co w przypadku exploita w bibliotece systemowej pozwala na przejęcie kontroli nad całym hostem.
}

\subsection{Bezpieczeństwo Zależności}
\finding{warning}{Użycie Bibliotek w Wersjach Beta}{Server/src/Terminal.Backend.Infrastructure/Terminal.Backend.Infrastructure.csproj}{
    Krytyczne komponenty OpenTelemetry (np. \texttt{OpenTelemetry.Instrumentation.EntityFrameworkCore}) są w wersji \texttt{1.0.0-beta.12}. Wersje beta nie gwarantują stabilności ani bezpieczeństwa, a ich użycie w systemie produkcyjnym w 2026 roku jest wysoce ryzykowne.
}

\subsection{Sekrety i Konfiguracja}
\finding{warning}{Wyciek Szczegółów Błędów DB}{ansible/deploy-playbook.yml}{
    Ustawienie \texttt{Include Error Detail=true} w środowisku produkcyjnym powoduje, że komunikaty o błędach bazy danych (np. przy naruszeniu klucza obcego) zawierają szczegóły struktury tabel i wartości pól, co ułatwia rozpoznanie systemu przez napastnika.
}

\newpage
\section{Wydajność i Optymalizacja}

\subsection{Frontend (React & Vite)}
\finding{warning}{Brak Lazy Loading dla Tras}{Client/src/App.tsx}{
    Wszystkie komponenty stron są importowane statycznie. Powoduje to, że paczka \texttt{index.js} jest niepotrzebnie duża, co drastycznie wydłuża czas pierwszego załadowania aplikacji na urządzeniach mobilnych.
}

\finding{warning}{Błędna Logika Offline (PWA)}{Client/src/App.tsx}{
    Aplikacja PWA powinna pozwalać na przeglądanie danych w trybie offline. Tymczasem w \texttt{App.tsx} wiele tras jest renderowanych warunkowo: \texttt{\{online \&\& ...\}}. Jeśli użytkownik straci zasięg, aplikacja przestanie renderować formularze, zamiast pozwolić na wpisanie danych i ich synchronizację po odzyskaniu połączenia.
}

\subsection{Backend (EF Core)}
\finding{warning}{N+1 i Nieoptymalne Zapytania}{Server/src/Terminal.Backend.Infrastructure/DAL/Handlers/Projects/GetProjectQueryHandler.cs}{
    Kod zawiera jawny komentarz \texttt{// FIXME: .Include(p => p.Processes)}. Zamiast jednego zapytania z JOIN, system wykonuje dwa osobne zapytania do bazy danych, co przy dużej skali projektów doprowadzi do degradacji wydajności.
}

\newpage
\section{Jakość Kodu i Procesy CI/CD}

\subsection{Testowanie (Największy Słaby Punkt)}
\finding{danger}{Krytyczny Brak Testów i Błąd w CI}{.github/workflows/terminal-server-ci.yml}{
    Pipeline CI dla backendu wykonuje jedynie \texttt{dotnet build}. Całkowicie pominięto krok \texttt{dotnet test}. W połączeniu z niemal zerową liczbą testów jednostkowych w katalogu \texttt{test/}, system jest całkowicie nieodporny na regresje.
}

\subsection{Standardy Programistyczne}
\successbox{Czysta Architektura}{
    Separacja warstw (\texttt{Core}, \texttt{Application}, \texttt{Infrastructure}, \texttt{Api}) jest wzorowa. Wykorzystanie wzorca MediatR pozwala na łatwe rozszerzanie logiki biznesowej bez modyfikacji istniejącego kodu.
}

\successbox{Obserwowalność (Observability)}{
    Integracja z HyperDX i OpenTelemetry zapewnia wgląd w działanie systemu na poziomie światowej klasy. To najmocniejszy punkt techniczny projektu.
}

\newpage
\section{Wnioski i Rekomendacje}

\subsection{Priorytet 1: Natychmiastowe (Security Fixes)}
\begin{enumerate}
    \item \textbf{CORS:} Ograniczyć \texttt{AllowAnyOrigin} do konkretnych domen.
    \item \textbf{Security Headers:} Wdrożyć middleware dodający nagłówki \texttt{HSTS, CSP, X-Frame-Options}.
    \item \textbf{Rate Limiting:} Dodać limitowanie zapytań dla endpointów logowania i zapisu.
\end{enumerate}

\subsection{Priorytet 2: Krótkoterminowe (Quality & Performance)}
\begin{enumerate}
    \item \textbf{CI/CD:} Dodać krok \texttt{dotnet test} do pipeline'u i zacząć pisać testy jednostkowe dla \texttt{Application}.
    \item \textbf{Docker:} Przejść na obrazy \texttt{alpine} lub \texttt{distroless} i zdefiniować nieuprzywilejowanego użytkownika.
    \item \textbf{React:} Wdrożyć \texttt{React.lazy()} i \texttt{Suspense} dla tras w \texttt{App.tsx}.
\end{enumerate}

\subsection{Priorytet 3: Średnioterminowe (Architecture)}
\begin{enumerate}
    \item \textbf{IDOR:} Zaimplementować \texttt{Resource-based Authorization}.
    \item \textbf{PWA:} Przebudować logikę offline, aby formularze były dostępne bez połączenia z siecią.
    \item \textbf{EF Core:} Usunąć problemy N+1 poprzez poprawne użycie \texttt{.Include()}.
\end{enumerate}

\vspace{2cm}
\begin{center}
    \textit{Koniec Raportu}
\end{center}

\end{document}
